\scanpage{51}
% archive sheet number: 50
Donc en admettant ce lemme, on trouverait que tt syst. $v_1\ldots v_n$ peut se modifier \unsure{galoisiennement|galois\textsuperscript{ment}} de façon à donner lieu à un pt \deleted{générique} de $V(\deleted{\unsure{\widehat{\mathbb Q}}}\widehat{\mathbb Z})$ \deleted{générique} tel que \uline{pour tt $l$}, \deleted{\unsure{$V_l$}} son image dans $V(( \mathbb Z_l)$ soit générique rel. à $\mathbb Q$,

D'ailleurs, en sens inverse, si on admet que $G$ satisfait à la condition d'approximation pour la place à l'infini, $G(\mathbb Q)$ dense \unsure{dans|ds} $G(\widehat{\mathbb Q})$,

\margin{Moyennant \inserted{le ``lemme'' et} la condition d'approximation sur $G$, la conjecture de Tate renforcée \unsure{i.e.|o.c.} \uline{équivaut} à \deleted{\unsure{la}} l'énoncé d'indépendance des réseaux $M_{v_i}^{\mathbb Z}$ dans $M^{\mathbb A_f}$ ($\mathbb A_f=\prod_{l\ \mathrm{fini}}\mathbb Z_l$)}
% slanted
\unsure{alors} donc la classe envisagée \unsure{$\Delta.a$|$U.a$} contient un pt rationnel sur $\mathbb Q$. Donc cela prouve que quitte à modifier $v_i$ \unsure{galoisiennement|galois\textsuperscript{ment}}, on trouve des réseaux $M_{v_i}^{\mathbb Z}$ tous commensurables.
[On fera cependant attention que le choix des $\sigma_i$ par lesquels on \inserted{doit} modifier $v_i$ \uline{dépend de $M$} ; il est possible que pour $M$ variable, une place $v$ soit bien déterminée par la connaissance du foncteur fibre $\unsure{\xi_v|\zeta_v}$].
